\section{Trabajos relacionados}
\label{sec:trabajos}

\subsection{Detección de discurso de odio y la herencia de HatEval}

La detección de discurso de odio es una tarea consolidada en el procesamiento
del lenguaje natural, con revisiones sistemáticas que cubren definiciones,
corpus y enfoques \cite{schmidt2017survey,fortuna2018survey,jahan2023systematic}.
Un hito para el español y el multilingüismo fue la tarea 5 de SemEval-2019
(HatEval), que introdujo la anotación conjunta de odio y carácter dirigido
contra inmigrantes y mujeres \cite{basile2019hateval}. Los sistemas
participantes combinaron rasgos lingüísticos, n-gramas de caracteres y modelos
lineales o neuronales \cite{argotavega2019mineriaunam,ninaalcocer2019haterecognizer,plazadelarco2019sinai}.
La lección que heredamos es doble: la estructura de las etiquetas importa y
los corpus son pequeños, ruidosos y sensibles al dominio.

En español, HaterNet mostró que un clasificador lineal sobre representaciones
dispersas podía competir con alternativas más complejas en un corpus de
Twitter \cite{pereirakohatsu2019haternet}; MexA3T exploró la agresividad y el
odio en variantes del español de México \cite{alvarezcarmona2018mexa3t,aragon2019mexa3t};
otros trabajos abordaron la misoginia y la combinación de rasgos lingüísticos
\cite{garcia2021misogyny,garcia2022featurecombination,frenda2018misogyny,arango2019hate}.
Los transformers han elevado el techo en corpus grandes
\cite{plazadelarco2021comparing,plazadelarco2021multitask}, pero su ventaja se
reduce cuando el corpus es pequeño, el dominio está normalizado y el
presupuesto de cómputo es limitado. La restricción de esta competencia
---sin modelos preentrenados--- vuelve relevante la pregunta inversa: cuánto
se puede extraer de una representación construida desde cero.

\subsection{Representaciones subléxicas y modelos lineales}

Los n-gramas de caracteres son una representación clásica para texto corto y
ruidoso. Cavnar y Trenkle los usaron para categorización de texto
\cite{cavnar1994ngram} y McNamee y Mayfield mostraron su eficacia en
recuperación multilingüe \cite{mcnamee2004character}. La representación
TF-IDF \cite{salton1988term} combinada con regresión logística o máquinas
de vectores sigue siendo una línea base difícil de superar en corpus
pequeños. Los optimizadores importan: \textsc{liblinear} resuelve el problema
primal con coordenadas \cite{fan2008liblinear}, el descenso estocástico
escala a matrices dispersas \cite{bottou2010large} y SAGA acelera la
regularización ElasticNet \cite{defazio2014saga,zou2005regularization}. Los
ensambles de modelos diversos ---bagging, stacking o mezclas de
probabilidades--- reducen varianza y suavizan los errores de frontera
\cite{breiman1996bagging,wolpert1992stacked}. Nuestra propuesta se sitúa en
esa tradición, pero añade una capa de decisión explícita y auditable.

\subsection{Umbrales, calibración y estructura de etiquetas}

En clasificación desbalanceada, el umbral de 0.5 no es óptimo para la $F_1$
\cite{japkowicz2002class,lipton2014thresholding}. La literatura ofrece
soluciones analíticas y numéricas para maximizar la $F_1$ por clase
\cite{nan2012optimizing,dembczynski2013optimizing,parambath2014optimizing,berger2020threshold}
y advierte que la métrica tiene una geometría propia: no todas las celdas de
la matriz de confusión pesan igual \cite{flach2003geometry,saito2015precision}.
La calibración de probabilidades es el paso previo: la temperatura y el
escalado de Platt corrigen el exceso de confianza de los modelos
\cite{guo2017calibration,platt1999probabilistic}. Finalmente, la estructura
jerárquica de las etiquetas ---TR solo tiene sentido si HS está activo---
puede explotarse con reglas ontológicas o con clasificadores jerárquicos
\cite{silla2011survey,fortuna2019hierarchically}. Nuestro pipeline combina
las tres ideas: calibración, umbrales desacoplados y una regla ascendente.

\subsection{Calidad de datos, fuga y comparación estadística}

Los corpus de odio contienen near-duplicados, retuits, hilos de un mismo
usuario y etiquetas ambiguas \cite{founta2018crowdsourcing,davidson2017automated,davidson2019racial,wiegand2019biased}.
Si la validación reparte copias casi idénticas entre entrenamiento y
validación, la métrica local se infla: es la fuga por agrupamiento descrita
por Kaufman et al. \cite{kaufman2012leakage}. La respuesta metodológica es
agrupar antes de partir y verificar que ningún grupo cruce la frontera
\cite{kohavi1995study,kaufman2012leakage}. La comparación entre modelos
tampoco debería apoyarse en una única diferencia de medias: el bootstrap
agrupado, las pruebas pareadas de McNemar y los tests de Dietterich y Demšar
son el estándar para decidir si una mejora es distinguishable del ruido
\cite{efron1979bootstrap,mcnemar1947note,dietterich1998approximate,demsar2006statistical}.
Ese es el marco que adoptamos.

\subsection{Posicionamiento}

Este trabajo no propone una arquitectura nueva. Propone una \emph{disciplina}
de construcción y evaluación: representación subléxica multiescala,
validación agrupada, umbrales por etiqueta, calibración y una capa de
decisión que preserva la prevalencia. La contribución diferencial es el
análisis formal de esa capa y la separación explícita entre lo que se valida
localmente y lo que se ajusta contra el leaderboard. Frente a los sistemas
basados en transformers, el nuestro es más simple y más barato; frente a las
líneas base clásicas, incorpora una ingeniería de decisión que resulta
decisiva en la competencia.
